\chapter{Extensões de corpos e teoria de Galois}\label{ch:b3:galois}

Toda equação pode ser resolvida por radicais, como sugerem a fórmula
quadrática e as fórmulas de Cardano para o \omterm{def:b3:galois:extension}{grau} três? É possível
trissectar um ângulo com régua e compasso? As duas perguntas, abertas
por séculos, recebem resposta --- negativa --- de uma única ideia de
Évariste Galois: associar a cada polinômio um \emph{grupo} finito de
simetrias de suas raízes e ler a resposta no grupo. Este capítulo
constrói o dicionário: extensões de corpos e \omterm{def:b3:galois:extension}{graus}, \omterm{thm:b3:galois:splitting}{corpos de
decomposição} e \omterm{def:b3:galois:closure}{fechos algébricos}, \omterm{thm:b3:galois:finitefields}{corpos finitos} (uma teoria completa
--- e a prometida ciclicidade de $\mathbb F_q^\times$), separabilidade e, então, a
própria \omterm{thm:b3:galois:fundamental}{correspondência de Galois}, com demonstrações integrais.
Colhemos: a impossibilidade das construções clássicas, a estrutura dos
corpos ciclotômicos e a insolubilidade da quíntica por radicais --- a
simplicidade de $A_5$, obtida no~\cref{ch:b3:groups}, atingindo seu alvo.

\section{Extensões, grau, algebricidade}

\begin{definition}\label{def:b3:galois:extension}
Uma \emph{extensão de corpos}\index{extensão de corpos} $L/K$ é um
corpo $L$ que contém $K$ como subcorpo; $L$ é então um
$K$-espaço vetorial, e o \emph{grau}\index{grau de uma extensão}
$[L:K]$ é sua dimensão. A extensão é \emph{finita} se $[L:K] <
\infty$. A
\emph{característica} de um corpo é o gerador $\geq 0$ do núcleo de
$\Z \to K$, $n \mapsto n\cdot 1$: ela vale $0$ ou um primo $p$; correspondentemente,
$K$ contém um menor subcorpo (\emph{corpo primo}) isomorfo a $\Q$
ou a $\mathbb
F_p = \Z/p\Z$.
\end{definition}

\begin{theorem}[Lei da torre]\label{thm:b3:galois:tower}
Se $K \subseteq L \subseteq M$, então $[M:K] = [M:L]\,[L:K]$: se $(e_i)$ é uma base de $L$ sobre $K$ e
$(f_j)$ é uma base de $M$ sobre $L$, então $(e_if_j)$ é uma base de
$M$ sobre $K$.\index{lei da torre}
\end{theorem}

\begin{proof}
Geração: $x \in M$ escreve-se $x = \sum_j \lambda_jf_j$ ($\lambda_j \in L$), e cada $\lambda_j = \sum_i \mu_{ij}e_i$ ($\mu_{ij} \in K$): $x = \sum_{i,j}\mu_{ij}e_if_j$.
Independência: $\sum_{i,j}\mu_{ij}e_if_j = 0$ reescreve-se como $\sum_j (\sum_i
\mu_{ij}e_i)f_j = 0$; as somas internas estão em
$L$, logo se anulam (os $f_j$ são independentes sobre $L$); e
então todos os $\mu_{ij} = 0$ (os $e_i$ são independentes sobre $K$).
\end{proof}

\begin{definition}\label{def:b3:galois:algebraic}
Sejam $L/K$ e $\alpha \in L$. Se algum $P \in K[X]$ não nulo satisfaz $P(\alpha) = 0$, dizemos
que $\alpha$ é \emph{algébrico}\index{elemento algébrico} sobre
$K$; o gerador mônico $\pi_\alpha$ do \omterm{def:b3:rings:ideal}{ideal} $\{P : P(\alpha) = 0\}$ de $K[X]$ é seu
\emph{polinômio minimal}\index{polinômio minimal (teoria de corpos)},
um polinômio \omterm{def:b3:rings:divisibility}{irredutível} ($\pi = QR$ com $Q(\alpha) = 0$ força $R$ constante, pela
minimalidade do \omterm{def:b3:galois:extension}{grau}). Caso contrário, $\alpha$ é
\emph{transcendente}. Escrevemos $K(\alpha)$ para o menor subcorpo de $L$
que contém $K$ e $\alpha$, e $K[\alpha]$ para o menor subanel.
\end{definition}

\begin{theorem}\label{thm:b3:galois:simple}
Se $\alpha$ é \omterm{def:b3:galois:algebraic}{algébrico} sobre $K$ com $d = \deg\pi_\alpha$, então
\[
K(\alpha) = K[\alpha] \cong K[X]/(\pi_\alpha),
\qquad [K(\alpha) : K] = d,
\]
com base $1, \alpha, \dots, \alpha^{d-1}$. Reciprocamente, se $[L:K] < \infty$, todo $\alpha \in L$ é \omterm{def:b3:galois:algebraic}{algébrico} de \omterm{def:b3:galois:extension}{grau}
$\deg \pi_\alpha$ que divide $[L:K]$.
\end{theorem}

\begin{proof}
A avaliação $K[X] \to L$, $P \mapsto P(\alpha)$, tem imagem $K[\alpha]$ e núcleo $(\pi_\alpha)$; como $\pi_\alpha$ é
\omterm{def:b3:rings:divisibility}{irredutível}, $K[X]/(\pi_\alpha)$ é um \emph{corpo} (\cref{prop:b3:rings:primemaximal}: no \omterm{def:b3:rings:pidufd}{DIP} $K[X]$,
um \omterm{def:b3:rings:divisibility}{irredutível} gera um \omterm{def:b3:rings:primemaximal}{ideal maximal}), de modo que $K[\alpha]$ é um corpo
contendo $K$ e $\alpha$: ele é igual a $K(\alpha)$. As classes de $1, X, \dots, X^{d-1}$
formam uma base do quociente (divisão euclidiana), donde a base e o
\omterm{def:b3:galois:extension}{grau}. Reciprocamente, se $[L:K]
= n < \infty$: os $1, \alpha, \dots, \alpha^n$ são dependentes, o que fornece
um polinômio aniquilador; então $[K(\alpha):K] =
\deg\pi_\alpha$ divide $n$ pela \omterm{thm:b3:galois:tower}{lei da torre}.
\end{proof}

\begin{corollary}\label{cor:b3:galois:algclosed}
Se $\alpha, \beta$ são \omterm{def:b3:galois:algebraic}{algébricos} sobre $K$, também o são $\alpha \pm
\beta$, $\alpha\beta$, $\alpha/\beta$
($\beta \ne 0$): os elementos de $L$ \omterm{def:b3:galois:algebraic}{algébricos} sobre $K$ formam um
subcorpo de $L$. Além disso, a algebricidade é transitiva: \omterm{def:b3:galois:algebraic}{algébrico}
sobre \omterm{def:b3:galois:algebraic}{algébrico} é \omterm{def:b3:galois:algebraic}{algébrico}.
\end{corollary}

\begin{proof}
$K(\alpha, \beta) = (K(\alpha))(\beta)$ é finita sobre $K(\alpha)$ ($\beta$ é \omterm{def:b3:galois:algebraic}{algébrico} sobre $K \subseteq K(\alpha)$) e
$K(\alpha)/K$ é finita: pela \omterm{thm:b3:galois:tower}{lei da torre}, $[K(\alpha,\beta):K] <
\infty$, e todo elemento de $K(\alpha, \beta)$ ---
inclusive os quatro listados --- é \omterm{def:b3:galois:algebraic}{algébrico} (\cref{thm:b3:galois:simple}).
Transitividade: se $\beta$ é \omterm{def:b3:galois:algebraic}{algébrico} sobre $L$ e $L/K$ é
\omterm{def:b3:galois:algebraic}{algébrica}, os coeficientes $c_0, \dots, c_{m-1}$ de $\pi_{\beta/L}$ geram uma extensão finita
$F = K(c_0, \dots,
c_{m-1})$ de $K$ (\omterm{thm:b3:galois:tower}{lei da torre} aplicada repetidamente), e $F(\beta)/F$ é
finita: $[F(\beta):K] < \infty$, logo $\beta$ é \omterm{def:b3:galois:algebraic}{algébrico} sobre $K$.
\end{proof}

\begin{example}\label{ex:b3:galois:degrees}
$[\Q(\sqrt2):\Q] = 2$, $[\Q(\sqrt[3]2):\Q] = 3$ ($X^3 - 2$ é \omterm{def:b3:rings:divisibility}{irredutível}, por \omterm{thm:b3:rings:criteria}{Eisenstein}), $[\Q(\zeta_p):\Q] = p - 1$ para $\zeta_p = \eu^{2\iu\pi/p}$
($\Phi_p$ é \omterm{def:b3:rings:divisibility}{irredutível}, \cref{ex:b3:rings:eisenstein}). A \omterm{thm:b3:galois:tower}{lei da torre} já é uma
arma: $\sqrt[3]2 \notin \Q(\sqrt2)$, pois $3 \nmid 2$.
\end{example}

\section{Corpos de decomposição; fecho algébrico}

\begin{theorem}[Corpos de decomposição]\label{thm:b3:galois:splitting}
Seja $P \in K[X]$ não constante. Existe um \emph{corpo de
decomposição}\index{corpo de decomposição} de $P$ sobre $K$:
uma extensão $L =
K(\alpha_1, \dots, \alpha_n)$ gerada por raízes de $P$ na qual $P$ se
decompõe em fatores lineares. Ele é único a menos de
$K$-isomorfismo, e $[L:K] \leq (\deg P)!$.
\end{theorem}

\begin{proof}
\emph{Existência}, por indução sobre $\deg P$: escolha um fator
\omterm{def:b3:rings:divisibility}{irredutível} $Q$ de $P$; o corpo $K_1 = K[X]/(Q)$ contém a raiz $\alpha_1 = \bar X$ de
$Q$ e, portanto, de $P$; escreva $P = (X -
\alpha_1)P_1$ sobre $K_1$ e aplique
a indução a $P_1$ sobre $K_1$; os \omterm{def:b3:galois:extension}{graus} se multiplicam dando no
máximo $n(n-1)\cdots = n!$.

A \emph{unicidade} decorre do \emph{lema de extensão de isomorfismos},
que é mais forte: seja $\sigma \colon K \to K'$ um isomorfismo, $P \in K[X]$, e seja $P^\sigma$ o
polinômio de coeficientes transportados; sejam $L, L'$ \omterm{thm:b3:galois:splitting}{corpos de
decomposição} de $P, P^\sigma$; então $\sigma$ se estende a um isomorfismo
$L \to L'$. Indução sobre $[L:K]$: se $P$ se decompõe em $K$,
então $L = K$ e $L' = K'$ ($P^\sigma$ se decompõe em $K'$, e $L'$ é
gerado por suas raízes). Caso contrário, escolha uma raiz $\alpha \in L \setminus K$ de um
fator \omterm{def:b3:rings:divisibility}{irredutível} $Q$ de $P$ com $\deg Q \geq 2$; $Q^\sigma$ é um fator
\omterm{def:b3:rings:divisibility}{irredutível} de $P^\sigma$, com uma raiz $\beta \in L'$; então
\[
K(\alpha) \cong K[X]/(Q) \xrightarrow{\ \sigma\ }
K'[X]/(Q^\sigma) \cong K'(\beta)
\]
estende $\sigma$ com $\alpha \mapsto \beta$. Ora, $L$ é um \omterm{thm:b3:galois:splitting}{corpo de decomposição}
de $P$ sobre $K(\alpha)$, e $L'$ o é de $P^\sigma$ sobre $K'(\beta)$, com $[L : K(\alpha)] < [L:K]$: a
indução estende ainda mais, até $L \to L'$.
\end{proof}

\begin{definition}\label{def:b3:galois:closure}
Um corpo $\Omega$ é \emph{algebricamente fechado}\index{corpo
algebricamente fechado} se todo polinômio não constante de $\Omega[X]$ tem
uma raiz em $\Omega$ (e portanto se decompõe). Um \emph{fecho
algébrico}\index{fecho algébrico} de $K$ é uma extensão \omterm{def:b3:galois:algebraic}{algébrica}
$\bar K/K$ com $\bar K$ \omterm{def:b3:galois:algebraic}{algebricamente} fechado.
\end{definition}

\begin{theorem}[Steinitz]\label{thm:b3:galois:steinitz}
Todo corpo $K$ tem um \omterm{def:b3:galois:closure}{fecho algébrico}, único a menos de
$K$-isomorfismo.
\end{theorem}

\begin{proof}
\emph{Existência (construção de Artin).} Seja $R = K[(X_f)_f]$ o anel de
polinômios com uma variável $X_f$ para cada $f \in K[X]$ mônico não
constante, e seja $I$ o \omterm{def:b3:rings:ideal}{ideal} gerado por todos os $f(X_f)$. O
\omterm{def:b3:rings:ideal}{ideal} $I$ é próprio: uma relação $1 = \sum_{i=1}^r g_i\,
f_i(X_{f_i})$ envolve um número finito de
polinômios; em um \omterm{thm:b3:galois:splitting}{corpo de decomposição} comum $E$ de $f_1\cdots f_r$, escolha
raízes $\alpha_i$ de $f_i$ e avalie $X_{f_i} \mapsto \alpha_i$ (as demais variáveis em
$\mapsto 0$): $1 = 0$, absurdo. Seja $\mathfrak m \supseteq I$ maximal (\cref{thm:b3:rings:krull}; Zorn)
e ponha $K_1 = R/\mathfrak
m$: uma \omterm{def:b3:galois:extension}{extensão de corpos} de $K$ em que todo $f \in
K[X]$ não
constante tem uma raiz, a saber $\bar X_f$, e que é \omterm{def:b3:galois:algebraic}{algébrica} sobre $K$
(ela é gerada pelos $\bar X_f$, cada um \omterm{def:b3:galois:algebraic}{algébrico}). Itere: $K \subseteq K_1 \subseteq K_2 \subseteq \cdots$, onde $K_{n+1}$
faz com $K_n$ o que $K_1$ fez com $K$, e ponha $\Omega =
\bigcup_n K_n$, um
corpo. Todo $g \in \Omega[X]$ não constante tem seus finitos coeficientes em algum
$K_n$; um fator \omterm{def:b3:rings:divisibility}{irredutível} de $g$ sobre $K_n$ tem uma raiz
em $K_{n+1} \subseteq
\Omega$: $\Omega$ é \omterm{def:b3:galois:closure}{algebricamente fechado}, e é \omterm{def:b3:galois:algebraic}{algébrico} sobre
$K$ (cada $K_n$ o é, por transitividade, \cref{cor:b3:galois:algclosed}):
$\Omega$ é um \omterm{def:b3:galois:closure}{fecho algébrico}.

\emph{Unicidade.} Sejam $\Omega, \Omega'$ dois \omterm{def:b3:galois:closure}{fechos algébricos}. Considere o
conjunto dos pares $(E, \tau)$ em que $K
\subseteq E \subseteq \Omega$ e $\tau\colon E \to \Omega'$ é um $K$-mergulho,
ordenado por extensão; ele é não vazio ($(K,
\mathrm{id})$) e indutivo (a reunião de
uma cadeia), de modo que Zorn fornece um par maximal $(E_0, \tau_0)$. Se $E_0 \neq \Omega$,
escolha $\alpha \in
\Omega\setminus E_0$: $\pi_{\alpha/E_0}$ é levado em um polinômio sobre $\tau_0(E_0)$ que tem uma
raiz $\beta$ no \omterm{def:b3:galois:closure}{corpo algebricamente fechado} $\Omega'$, e $\tau_0$
se estende a $E_0(\alpha) \to
\Omega'$ ($\alpha \mapsto \beta$), contradizendo a maximalidade. Existe,
portanto, um $K$-mergulho $\tau \colon \Omega \to \Omega'$; sua imagem, isomorfa a
$\Omega$, é \omterm{def:b3:galois:algebraic}{algebricamente} fechada, e $\Omega'$ é \omterm{def:b3:galois:algebraic}{algébrico} sobre ela:
para $x \in \Omega'$, $\pi_{x/\tau(\Omega)}$ se decompõe sobre $\tau(\Omega)$, logo $x \in
\tau(\Omega)$. Assim
$\tau$ é sobrejetor: é um isomorfismo.
\end{proof}

\begin{remark}\label{rem:b3:galois:closureexamples}
Para $K = \Q$ é possível evitar o maquinário transfinito: os
números \omterm{def:b3:galois:algebraic}{algébricos} $\bar\Q = \{z \in \C : z \text{ algebraic over
} \Q\}$ formam um \omterm{def:b3:galois:closure}{fecho algébrico} --- um subcorpo de
$\C$ pelo~\cref{cor:b3:galois:algclosed}, \omterm{def:b3:galois:closure}{algebricamente fechado} porque $\C$ o
é (d'Alembert--Gauss, demonstrado por análise complexa no~\cref{ch:b3:holomorphic}), e as raízes de polinômios sobre $\bar\Q$ são
\omterm{def:b3:galois:algebraic}{algébricas} sobre $\Q$ por transitividade.
\end{remark}

\section{Corpos finitos}

\begin{theorem}\label{thm:b3:galois:finitefields}
Sejam $p$ primo, $n \geq 1$, $q = p^n$.
\begin{enumerate}
  \item Um \omterm{thm:b3:galois:finitefields}{corpo finito} tem por cardinalidade uma potência de primo e,
        para cada $q$, existe exatamente um corpo $\mathbb F_q$ com
        $q$ elementos a menos de isomorfismo: o \omterm{thm:b3:galois:splitting}{corpo de
        decomposição} de $X^q - X$ sobre $\mathbb F_p$.\index{corpo finito}
  \item O \emph{Frobenius}\index{endomorfismo de Frobenius} $F
        \colon x \mapsto x^p$ é um
        automorfismo de $\mathbb
        F_q$, e o grupo de automorfismos de $\mathbb F_q$ é
        cíclico de ordem $n$, gerado por $F$.
  \item $\mathbb F_{p^m}$ mergulha em $\mathbb F_{p^n}$ se, e somente se, $m \mid
        n$.
\end{enumerate}
\end{theorem}

\begin{proof}
(1) Um \omterm{thm:b3:galois:finitefields}{corpo finito} $E$ tem característica $p > 0$ e é um
$\mathbb F_p$-espaço vetorial de dimensão finita: $\abs E = p^n$. Seu grupo
multiplicativo tem ordem $q - 1$, de modo que todo $x \in E$ satisfaz
$x^q = x$: $E$ é formado por $q$ raízes de $X^q - X$ e, portanto, é um
\omterm{thm:b3:galois:splitting}{corpo de decomposição} desse polinômio sobre $\mathbb F_p$ --- o que determina
$E$ a menos de isomorfismo (\cref{thm:b3:galois:splitting}). Reciprocamente, em um
\omterm{thm:b3:galois:splitting}{corpo de decomposição} $L$ de $X^q - X$, o conjunto $E$ de suas
raízes é um \emph{subcorpo}: $(x + y)^q = x^q + y^q$ iterando o sonho do calouro $(a+b)^p = a^p + b^p$
($p \mid \binom pk$), e $(xy)^q = x^qy^q$, $(x^{-1})^q = (x^q)^{-1}$; ele tem exatamente $q$ elementos, pois
$X^q - X$ é \omterm{def:b3:galois:separable}{separável}: sua derivada é $qX^{q-1} - 1 = -1$ (já que $p \mid q$), coprima com
ele, de modo que não há raízes múltiplas. Assim $L = E$ tem $q$
elementos.

(2) $F$ é um morfismo de corpos (sonho do calouro), injetor (são
corpos) e, portanto, bijetor no conjunto finito $\mathbb F_q$. Tem-se $F^n =
\mathrm{id}$
($x^q = x$), e nenhuma potência menor é a identidade: $F^m = \mathrm{id}$ significaria
que todos os $q$ elementos são raízes de $X^{p^m} - X$, o que força $p^m \geq q$.
Logo $\langle F\rangle$ é cíclico de ordem $n$; e não há outros automorfismos,
pela cota $\abs{\operatorname{Aut}} \leq [\,\mathbb F_q :
\mathbb F_p\,] = n$ demonstrada adiante (\cref{prop:b3:galois:embeddings} com $L = \mathbb F_q$, $K =
\mathbb F_p$:
os automorfismos fixam o corpo primo).

(3) Se $\mathbb F_{p^m} \subseteq \mathbb F_{p^n}$, a \omterm{thm:b3:galois:tower}{lei da torre} dá $p^n = (p^m)^d$: $m \mid n$. Reciprocamente, se $m \mid n$,
então $p^m - 1 \mid p^n - 1$ (soma geométrica), de modo que $X^{p^m} - X$ divide $X^{p^n} - X$ (mesmo
argumento nos expoentes: $X^{a} - 1 \mid X^{b}
- 1$ quando $a \mid b$), e as raízes do primeiro
dentro de $\mathbb F_{p^n}$ formam o subcorpo procurado, de cardinalidade
$p^m$ (separabilidade como em (1)).
\end{proof}

\begin{theorem}[Ciclicidade]\label{thm:b3:galois:cyclic}
Todo subgrupo finito do grupo multiplicativo de um corpo é cíclico. Em
particular, $\mathbb F_q^\times \cong
\Z/(q-1)\Z$.\index{grupo multiplicativo de um corpo finito}
\end{theorem}

\begin{proof}
Seja $G \leq K^\times$ finito. Pelo teorema de estrutura (\cref{cor:b3:modules:abelian}), $G \cong \Z/d_1 \times \dots
\times \Z/d_s$
com $d_1 \mid \dots \mid d_s$. Todo $x \in G$ satisfaz então $x^{d_s} = 1$; mas $X^{d_s} - 1$ tem no máximo
$d_s$ raízes no corpo $K$: $\abs G = d_1\cdots d_s \leq d_s$, o que força $s = 1$:
$G$ é cíclico.
\end{proof}

\begin{example}\label{ex:b3:galois:f8}
$\mathbb F_8 = \mathbb F_2[X]/(X^3 + X + 1)$: a cúbica não tem raiz em $\mathbb F_2$ e, portanto, é \omterm{def:b3:rings:divisibility}{irredutível}.
Escrevendo $\omega =
\bar X$: $\mathbb F_8^\times$ é cíclico de ordem $7$, de modo que
\emph{todo} elemento $\neq 0, 1$ é gerador. Os subcorpos de $\mathbb F_{p^{12}}$ formam o
reticulado dos divisores de $12$: $\mathbb F_p, \mathbb F_{p^2}, \mathbb F_{p^3}, \mathbb F_{p^4},
\mathbb F_{p^6}, \mathbb F_{p^{12}}$ --- um primeiro exemplo
completo da \omterm{thm:b3:galois:fundamental}{correspondência de Galois}.
\end{example}

\section{Separabilidade e mergulhos}

\begin{definition}\label{def:b3:galois:separable}
Um polinômio $P \in K[X]$ é \emph{separável}\index{polinômio separável} se não
tem raiz múltipla em um \omterm{thm:b3:galois:splitting}{corpo de decomposição} --- equivalentemente,
$\gcd(P, P') = 1$ (uma raiz múltipla é uma raiz comum; reciprocamente, sobre o
\omterm{thm:b3:galois:splitting}{corpo de decomposição}, uma raiz comum é múltipla; e o mdc não muda por
extensão de corpo, pelo argumento do~\cref{cor:b3:modules:descent}). Um \omterm{def:b3:galois:algebraic}{elemento
algébrico} é separável se seu \omterm{def:b3:galois:algebraic}{polinômio minimal} o for; uma extensão
$L/K$ é separável se todos os seus elementos o forem.
\end{definition}

\begin{proposition}\label{prop:b3:galois:perfect}
Um $P \in K[X]$ \emph{\omterm{def:b3:rings:divisibility}{irredutível}} é \omterm{def:b3:galois:separable}{separável}, exceto se $P' = 0$, o que
força $\operatorname{char} K = p > 0$ e $P \in K[X^p]$. Consequentemente, toda extensão \omterm{def:b3:galois:algebraic}{algébrica} de um
corpo de característica $0$, e de um \omterm{thm:b3:galois:finitefields}{corpo finito}, é \omterm{def:b3:galois:separable}{separável}
(tais corpos são chamados \emph{perfeitos}\index{corpo perfeito}).
\end{proposition}

\begin{proof}
$\gcd(P, P')$ divide $P$; se não for $1$, a irredutibilidade força
$\gcd = P$ (a menos de uma constante), de modo que $P \mid P'$ com $\deg
P' < \deg P$:
$P' = 0$. Escrevendo $P = \sum a_kX^k$: $ka_k = 0$ para todo $k$, logo em
característica $0$ o polinômio $P$ é constante (excluído); em
característica $p$, $a_k = 0$ a menos que $p \mid k$: $P =
Q(X^p)$. Sobre um \omterm{thm:b3:galois:finitefields}{corpo
finito}, todo elemento é uma $p$-ésima potência (o \omterm{thm:b3:galois:finitefields}{Frobenius} é
sobrejetor), de modo que $Q(X^p) = \sum b_k^p X^{pk} = (\sum
b_kX^k)^p$ não é \omterm{def:b3:rings:divisibility}{irredutível}: ali tampouco pode
ocorrer $P' = 0$ para $P$ \omterm{def:b3:rings:divisibility}{irredutível}.
\end{proof}

\begin{proposition}[Contagem de mergulhos]\label{prop:b3:galois:embeddings}
Sejam $L = K(\alpha_1, \dots, \alpha_r)$ finita sobre $K$ e $\sigma \colon K \to \Omega$ um mergulho em um \omterm{def:b3:galois:closure}{corpo
algebricamente fechado}. Então o número de extensões de $\sigma$ a
$L$ é no máximo $[L:K]$, com igualdade se $L/K$ é \omterm{def:b3:galois:separable}{separável}.
Em particular, $\abs{\operatorname{Aut}_K(L)} \leq [L:K]$.
\end{proposition}

\begin{proof}
Indução sobre $[L:K]$ por etapas simples. Para $L = K(\alpha)$: uma extensão
$\tau$ fica determinada por $\tau(\alpha)$, que deve ser uma raiz em
$\Omega$ de $\pi_\alpha^\sigma$; reciprocamente, cada uma dessas raízes dá uma
extensão ($K(\alpha) \cong K[X]/(\pi_\alpha)$). O número de extensões é o número de raízes
\emph{distintas} de $\pi_\alpha^\sigma$ em $\Omega$: no máximo $\deg\pi_\alpha =
[K(\alpha):K]$, com
igualdade se, e somente se, $\pi_\alpha$ é \omterm{def:b3:galois:separable}{separável} (a separabilidade de
$\pi^\sigma$ e a de $\pi$ coincidem: o mdc com a derivada é preservado por
$\sigma$). No caso geral, fatore $L =
K(\alpha_1)(\alpha_2, \dots)$: as extensões de
$\sigma$ a $K(\alpha_1)$ são em número de $\leq [K(\alpha_1):K]$, e cada uma se estende de
$\leq [L : K(\alpha_1)]$ maneiras, por indução; multiplique (\omterm{thm:b3:galois:tower}{lei da torre}). No caso
\omterm{def:b3:galois:separable}{separável}, ambas as contagens são igualdades: os polinômios minimais
sobre o corpo maior $K(\alpha_1)$ dividem os de sobre $K$ e, portanto,
permanecem separáveis.
\end{proof}

\begin{theorem}[Elemento primitivo]\label{thm:b3:galois:primitive}
Toda extensão finita \emph{\omterm{def:b3:galois:separable}{separável}} é simples: $L =
K(\gamma)$ para algum
$\gamma$.\index{teorema do elemento primitivo}
\end{theorem}

\begin{proof}
Se $K$ é finito, $L$ também é, e um gerador $\gamma$ do
grupo cíclico $L^\times$ (\cref{thm:b3:galois:cyclic}) resolve. Seja $K$ infinito;
por indução, basta tratar $L =
K(\alpha, \beta)$. Seja $n = [L:K]$; pela~\cref{prop:b3:galois:embeddings} há
$n$ $K$-mergulhos distintos $\sigma_1, \dots, \sigma_n \colon L \to \Omega$ (com $\Omega$ um \omterm{def:b3:galois:closure}{fecho
algébrico}). O polinômio
\[
D(T)=\prod_{i<j}\bigl[\bigl(\sigma_i(\alpha)-\sigma_j(\alpha)
\bigr) + T\bigl(\sigma_i(\beta) - \sigma_j(\beta)\bigr)\bigr]
\]
não é identicamente nulo: um fator se anula identicamente somente se
$\sigma_i, \sigma_j$ coincidem tanto em $\alpha$ quanto em $\beta$ e,
portanto, em $L$ --- excluído para $i \neq j$. Como $K$ é infinito,
escolha $c \in
K$ com $D(c) \ne 0$: então os $n$ elementos $\sigma_i(\alpha +
c\beta)$ são dois a
dois distintos, de modo que $\gamma = \alpha + c\beta$ tem ao menos $n$ conjugados
distintos em $\Omega$, isto é, $\deg \pi_\gamma \geq n$: $[K(\gamma):K] \geq n = [L:K]$ força $L = K(\gamma)$.
\end{proof}

\section{A correspondência de Galois}

\begin{definition}\label{def:b3:galois:galois}
Uma extensão finita $L/K$ é \emph{de Galois}\index{extensão de
Galois} se é o \omterm{thm:b3:galois:splitting}{corpo de decomposição} de um polinômio \emph{\omterm{def:b3:galois:separable}{separável}}
sobre $K$. Seu \emph{grupo de Galois}\index{grupo de Galois} é
$\operatorname{Gal}(L/K) = \operatorname{Aut}_K(L)$, o grupo dos automorfismos do corpo $L$ que fixam $K$
ponto a ponto.
\end{definition}

\begin{proposition}\label{prop:b3:galois:galoisorder}
Se $L/K$ é \omterm{def:b3:galois:galois}{de Galois}, então $\abs{\operatorname{Gal}(L/K)} =
[L:K]$; além disso, $L/F$ é \omterm{def:b3:galois:galois}{de
Galois} para todo corpo intermediário $K
\subseteq F \subseteq L$, e todo $F$-mergulho
$L \to \Omega
\supseteq L$ tem imagem $L$ (\emph{normalidade}).
\end{proposition}

\begin{proof}
Seja $L$ o \omterm{thm:b3:galois:splitting}{corpo de decomposição} do \omterm{def:b3:galois:separable}{polinômio separável} $P$
sobre $K$, e fixe um \omterm{def:b3:galois:closure}{fecho algébrico} $\Omega \supseteq L$. A extensão $L/K$ é
\omterm{def:b3:galois:separable}{separável}: ela é gerada por raízes de $P$; a separabilidade de cada
elemento decorre do caso de igualdade abaixo, mas argumentemos
diretamente --- a~\cref{prop:b3:galois:embeddings} aplicada aos geradores (raízes do
\omterm{def:b3:galois:separable}{polinômio separável} $P$, cujos polinômios minimais dividem $P$)
fornece exatamente $[L:K]$ extensões de $K \hookrightarrow \Omega$ (na etapa indutiva, o
\omterm{def:b3:galois:algebraic}{polinômio minimal} de uma raiz de $P$ sobre um corpo intermediário
ainda divide $P$ e, portanto, é \omterm{def:b3:galois:separable}{separável}). Cada um desses
mergulhos $\tau \colon L \to \Omega$ permuta as raízes de $P$ ($\tau$ fixa os
coeficientes), e $L$ é gerado por elas: $\tau(L) = L$. Logo mergulhos
$=$ automorfismos: $\abs{\operatorname{Gal}(L/K)} = [L:K]$. Para $F$ intermediário: $L$ é
também o \omterm{thm:b3:galois:splitting}{corpo de decomposição} de $P$ sobre $F$, e $P$
continua \omterm{def:b3:galois:separable}{separável}: $L/F$ é \omterm{def:b3:galois:galois}{de Galois}; o mesmo argumento dá a
normalidade sobre $F$.
\end{proof}

\begin{lemma}[Artin]\label{lem:b3:galois:artin}
Sejam $G$ um grupo finito de automorfismos de um corpo $L$ e
$K = L^G = \{x : \sigma(x) = x\ \forall\sigma \in G\}$ seu corpo fixo. Então $[L : L^G] \leq \abs G$.
\end{lemma}

\begin{proof}
Sejam $n = \abs G$, $G = \{\sigma_1, \dots, \sigma_n\}$, e suponha que $x_1, \dots, x_{n+1} \in L$ sejam linearmente
independentes sobre $K$. O sistema linear homogêneo de $n$
equações nas $n+1$ incógnitas $(c_j)$ sobre $L$,
\[
\sum_{j=1}^{n+1} c_j\,\sigma_i(x_j) = 0
\qquad (i = 1, \dots, n),
\]
tem solução não nula; escolha uma com o \emph{menor número} de entradas
não nulas, digamos $c_1, \dots, c_r \neq 0$ (renumerando), $r \geq 2$ (um único $c_j\sigma_i(x_j) = 0$ é
impossível), normalizada por $c_r = 1$. Nem todos os $c_j$ estão em
$K$: a equação relativa a $\sigma_i =
\mathrm{id}$ contradiria a independência;
digamos $c_1 \notin K$, de modo que $\tau(c_1) \ne c_1$ para algum $\tau \in G$. Aplique $\tau$ a
todas as equações: como $\tau\sigma_i$ percorre $G$, o vetor $(\tau(c_j))_j$ é outra
solução; subtraindo, $(c_j -
\tau(c_j))_j$ é uma solução com menos entradas não nulas
(a $r$-ésima entrada, $1 - 1 = 0$, se anula, e a primeira não) e não
nula: contradição. Portanto, quaisquer $n+1$ elementos são
dependentes: $[L:K]
\leq n$.
\end{proof}

\begin{theorem}[Teorema fundamental da teoria de Galois]
\label{thm:b3:galois:fundamental}
Seja $L/K$ uma \omterm{def:b3:galois:galois}{extensão de Galois} com grupo
$G =
\operatorname{Gal}(L/K)$.\index{correspondência de Galois}
\begin{enumerate}
  \item $L^G = K$.
  \item As aplicações $H \mapsto L^H$ e $F \mapsto
        \operatorname{Gal}(L/F)$ são bijeções mutuamente inversas e
        que invertem inclusões, entre os subgrupos de $G$ e os
        corpos intermediários $K \subseteq F \subseteq L$; além disso, $[L : L^H] = \abs H$ e $[L^H : K] = [G : H]$.
  \item $H \trianglelefteq G$ se, e somente se, $L^H/K$ é \omterm{def:b3:galois:galois}{de Galois}, e então a
        restrição induz $\operatorname{Gal}(L^H/K) \cong
        G/H$.
\end{enumerate}
\end{theorem}

\begin{proof}
(1) É claro que $K \subseteq L^G$. Reciprocamente, seja $\alpha \in L
\setminus K$; exibimos $\sigma \in G$ com
$\sigma(\alpha) \ne
\alpha$. O \omterm{def:b3:galois:algebraic}{polinômio minimal} $\pi_\alpha$ sobre $K$ tem \omterm{def:b3:galois:extension}{grau} $\geq 2$ e
é \omterm{def:b3:galois:separable}{separável} ($L/K$ é \omterm{def:b3:galois:separable}{separável}, \cref{prop:b3:galois:galoisorder}), de modo que tem
outra raiz $\beta \neq \alpha$ em um \omterm{def:b3:galois:closure}{fecho algébrico} $\Omega \supseteq L$. Estenda o $K$-mergulho
$K(\alpha) \to \Omega$, $\alpha \mapsto
\beta$, a um mergulho $\tau\colon L \to \Omega$ (\cref{prop:b3:galois:embeddings}); por normalidade
(\cref{prop:b3:galois:galoisorder}), $\tau(L) = L$, logo $\tau \in
G$, $\beta = \tau(\alpha) \in L$ e $\tau(\alpha) \neq
\alpha$.

(2) Para um subgrupo $H$: $L/L^H$ é \omterm{def:b3:galois:galois}{de Galois}
(\cref{prop:b3:galois:galoisorder}), e $\operatorname{Gal}(L/L^H) \supseteq H$ trivialmente, logo $[L:L^H] =
\abs{\operatorname{Gal}(L/L^H)} \geq \abs H$; o lema de Artin dá
$[L:L^H] \leq \abs H$: há igualdade, e $\operatorname{Gal}(L/L^H) = H$. Para um corpo intermediário $F$:
sendo $L/F$ \omterm{def:b3:galois:galois}{de Galois}, (1) aplicado a $L/F$ dá $L^{\operatorname{Gal}(L/F)} = F$. As duas
aplicações são mutuamente inversas; que invertem inclusões é evidente.
\omterm{def:b3:galois:extension}{Graus}: $[L:L^H] = \abs H$ acabou de ser demonstrado, e $[L^H:K] = [L:K]/[L:L^H] = \abs G/\abs H$.

(3) Para $\sigma \in G$ e $H \leq G$: $\sigma(L^H) =
L^{\sigma H\sigma^{-1}}$ (verificação direta). Pela bijeção, $\sigma(L^H) = L^H$
para todo $\sigma$ se, e somente se, $H \trianglelefteq G$. Ora, se $H \trianglelefteq G$, ponha
$F = L^H$: todo $\sigma \in
G$ se restringe a um automorfismo de $F$, o que dá um
morfismo $\rho
\colon G \to \operatorname{Aut}_K(F)$ de núcleo $\{\sigma :
\sigma\restriction_F = \mathrm{id}\} = \operatorname{Gal}(L/F) =
H$. Logo $G/H$ mergulha em $\operatorname{Aut}_K(F)$,
donde $\abs{\operatorname{Aut}_K(F)} \geq [G:H] = [F:K]$; a desigualdade inversa vale sempre (\cref{prop:b3:galois:embeddings}):
$\abs{\operatorname{Aut}_K(F)} = [F:K]$ e $\rho$ é sobrejetor. Resta ver que $F/K$ é \omterm{def:b3:galois:galois}{de Galois}:
$F$ é \omterm{def:b3:galois:separable}{separável} sobre $K$ (por estar dentro da extensão
\omterm{def:b3:galois:separable}{separável} $L/K$) e $F = K(\gamma)$ (\cref{thm:b3:galois:primitive}); o polinômio $\prod_{\sigma \in G/H}\bigl(X - \sigma(\gamma)\bigr)$
(produto sobre as imagens distintas, que estão em $F$: $\sigma(F) =
L^{\sigma H\sigma^{-1}} = L^H = F$ pela
normalidade de $H$) tem coeficientes fixados por $G$ e,
portanto, em $K$, por (1): trata-se de um \omterm{def:b3:galois:separable}{polinômio separável} de
$K[X]$ decomposto por $F$, cujas raízes geram $F$: $F/K$
é \omterm{def:b3:galois:galois}{de Galois}. Reciprocamente, se $F = L^H$ com $F/K$ \omterm{def:b3:galois:galois}{de Galois}, a
normalidade de $F$ (\cref{prop:b3:galois:galoisorder}, aplicada aos mergulhos $F \to \Omega$
restritos de elementos de $G$) dá $\sigma(F) = F$ para todo $\sigma \in G$, isto é,
$H
\trianglelefteq G$.
\end{proof}

\begin{omfigure}
\begin{tikzpicture}[xscale=1.5, yscale=1.15]
  \tikzset{fd/.style={font=\small}}
  % subgroup lattice of S3 (left, upside down to mirror fields)
  \node[fd] (G)  at (-3.2, 0)   {$S_3$};
  \node[fd] (A)  at (-4.4, 1.2) {$\langle(1\,2\,3)\rangle$};
  \node[fd] (T1) at (-3.6, 1.2) {$\langle(1\,2)\rangle$};
  \node[fd] (T2) at (-2.8, 1.2) {$\langle(1\,3)\rangle$};
  \node[fd] (T3) at (-2.0, 1.2) {$\langle(2\,3)\rangle$};
  \node[fd] (E)  at (-3.2, 2.4) {$\{e\}$};
  \draw (G) -- (A); \draw (G) -- (T1); \draw (G) -- (T2);
  \draw (G) -- (T3);
  \draw (E) -- (A); \draw (E) -- (T1); \draw (E) -- (T2);
  \draw (E) -- (T3);
  % field lattice (right)
  \node[fd] (Q)  at (1.6, 0)   {$\Q$};
  \node[fd] (W)  at (0.4, 1.2) {$\Q(\iu\sqrt3)$};
  \node[fd] (R1) at (1.4, 1.2) {$\Q(\sqrt[3]2)$};
  \node[fd] (R2) at (2.4, 1.2) {$\Q(j\sqrt[3]2)$};
  \node[fd] (R3) at (3.5, 1.2) {$\Q(j^2\sqrt[3]2)$};
  \node[fd] (L)  at (1.6, 2.4) {$L = \Q(\sqrt[3]2, j)$};
  \draw (Q) -- (W); \draw (Q) -- (R1); \draw (Q) -- (R2);
  \draw (Q) -- (R3);
  \draw (L) -- (W); \draw (L) -- (R1); \draw (L) -- (R2);
  \draw (L) -- (R3);
  \draw[->, omDef, thick, shorten >=6pt, shorten <=6pt]
    (-1.4, 1.2) -- (-0.6, 1.2);
\end{tikzpicture}

{\small A \omterm{thm:b3:galois:fundamental}{correspondência de Galois} para o \omterm{thm:b3:galois:splitting}{corpo de decomposição}
$L$ de $X^3 - 2$ sobre $\Q$ ($j = \eu^{2\iu\pi/3}$): os subgrupos de $\operatorname{Gal}(L/\Q) \cong S_3$ (à
esquerda, em ordem invertida) correspondem aos corpos intermediários (à
direita). O único \omterm{def:b3:groups:normal}{subgrupo normal} próprio $\langle(1\,2\,3)\rangle$ corresponde à única
subextensão $\Q(\iu\sqrt3)/\Q$ que é \omterm{def:b3:galois:galois}{de Galois}; os três subgrupos conjugados $\langle(i\,j)\rangle$
correspondem aos três corpos cúbicos conjugados $\Q(j^k\sqrt[3]2)$, nenhum deles
normal sobre $\Q$.}
\end{omfigure}

\section{Extensões ciclotômicas}

\begin{definition}\label{def:b3:galois:cyclotomic}
Sejam $n \geq 1$ e $\zeta_n = \eu^{2\iu\pi/n}$. O $n$-ésimo \emph{polinômio
ciclotômico}\index{polinômio ciclotômico} é $\Phi_n = \prod_{\gcd(k,n)=1,\ 1 \le k \le n} \bigl(X -
\zeta_n^k\bigr)$, de \omterm{def:b3:galois:extension}{grau} $\varphi(n)$;
agrupando as raízes de $X^n - 1$ por ordem exata, $X^n - 1 = \prod_{d \mid n}\Phi_d$, o que mostra por
indução que $\Phi_n \in \Z[X]$ (divisão euclidiana de polinômios inteiros mônicos).
\end{definition}

\begin{theorem}\label{thm:b3:galois:cyclotomicirred}
$\Phi_n$ é \omterm{def:b3:rings:divisibility}{irredutível} sobre $\Q$; logo $[\Q(\zeta_n) : \Q] =
\varphi(n)$ e
\[
\operatorname{Gal}\bigl(\Q(\zeta_n)/\Q\bigr) \;\cong\;
(\Z/n\Z)^\times,
\qquad \sigma_a(\zeta_n) = \zeta_n^a .
\]
A extensão $\Q(\zeta_n)/\Q$ é, portanto, \omterm{def:b3:galois:galois}{de Galois} com grupo \emph{abeliano}.
\end{theorem}

\begin{proof}
Seja $f = \pi_{\zeta_n}$, de modo que $\Phi_n = fg$ com $f, g \in \Z[X]$ mônicos (lema de Gauss,
\cref{lem:b3:rings:gauss}: os conteúdos se multiplicam e todos os polinômios são
mônicos). \emph{Afirmação: se $\zeta$ é raiz de $f$ e $p \nmid n$ é
primo, então $\zeta^p$ é raiz de $f$.} Caso contrário, $\zeta^p$ é raiz de
$g$ (é uma raiz $n$-ésima primitiva da unidade), de modo que
$\zeta$ é raiz de $g(X^p)$ e $f
\mid g(X^p)$ em $\Z[X]$ (\omterm{def:b3:galois:algebraic}{polinômio
minimal} e, de novo, Gauss). Reduza módulo $p$: $\bar g(X^p) = \bar g(X)^p$ (\omterm{thm:b3:galois:finitefields}{Frobenius} em
$\mathbb F_p[X]$: coeficiente a coeficiente $a^p = a$, mais o sonho do calouro), logo
$\bar f \mid \bar g^{\,p}$: $\bar f$ e $\bar g$ compartilham um fator \omterm{def:b3:rings:divisibility}{irredutível}, e
$\bar\Phi_n = \bar f\bar g$ tem um fator repetido. Então $X^n - \bar 1$ também tem; mas sua derivada
$\bar nX^{n-1}$ é coprima com ele ($p \nmid n$, e $0$ não é raiz): contradição.

Toda raiz primitiva $\zeta_n^k$ ($\gcd(k, n) = 1$) é obtida a partir de $\zeta_n$ por
potências sucessivas de primos que não dividem $n$ (fatore $k$):
a afirmação se propaga, de modo que toda raiz primitiva é raiz de
$f$: $f = \Phi_n$, \omterm{def:b3:rings:divisibility}{irredutível}. Consequentemente, $[\Q(\zeta_n):\Q] = \varphi(n)$, e $\Q(\zeta_n)$ é o \omterm{thm:b3:galois:splitting}{corpo
de decomposição} do \omterm{def:b3:galois:separable}{polinômio separável} $X^n - 1$ (todas as raízes são
potências de $\zeta_n$): é \omterm{def:b3:galois:galois}{de Galois}. Um automorfismo $\sigma$ leva
$\zeta_n$ em outra raiz primitiva $\zeta_n^{a(\sigma)}$, e $\sigma
\mapsto a(\sigma)$ é um morfismo injetor em
$(\Z/n\Z)^\times$; ambos os grupos têm ordem $\varphi(n)$: é um isomorfismo.
\end{proof}

\section{Régua e compasso}

\begin{definition}\label{def:b3:galois:constructible}
Identifique o plano com $\C$; parta de $\{0, 1\}$. Um ponto é
\emph{construtível}\index{número construtível} se pode ser obtido por
um número finito de interseções de retas que passam por dois pontos já
construídos e de círculos centrados em um ponto construído com raio
igual à distância entre dois pontos construídos.
\end{definition}

\begin{theorem}[Wantzel]\label{thm:b3:galois:wantzel}
$z \in \C$ é \omterm{def:b3:galois:constructible}{construtível} se, e somente se, existe uma torre $\Q = F_0
\subseteq F_1 \subseteq \dots \subseteq F_r$ com $[F_{i+1} :
F_i] = 2$
e $z \in F_r$. Em particular, um \omterm{def:b3:galois:constructible}{número construtível} é \omterm{def:b3:galois:algebraic}{algébrico} de \omterm{def:b3:galois:extension}{grau} uma
potência de $2$ sobre $\Q$.\index{teorema de Wantzel}
\end{theorem}

\begin{proof}
($\Rightarrow$) As coordenadas da interseção de duas retas que passam por
pontos de coordenadas em um subcorpo $F \subseteq \R$ resolvem um sistema linear
sobre $F$: elas permanecem em $F$. As interseções reta--círculo
e círculo--círculo levam, depois de eliminada a parte linear
(subtrair as duas equações dos círculos dá uma reta), a uma equação
quadrática sobre $F$: as novas coordenadas estão em $F$ ou em
$F(\sqrt d)$ para algum $d \in F$, $d > 0$. Por indução, todo ponto construído
tem coordenadas em uma torre de extensões quadráticas de $\Q$; e
$z = x + \iu y$ também está em uma torre quadrática (adjunte $\iu$: mais uma
etapa quadrática). A consequência sobre o \omterm{def:b3:galois:extension}{grau}: $[\Q(z):\Q]$ divide $[F_r : \Q] = 2^r$
(\omterm{thm:b3:galois:tower}{lei da torre}).

($\Leftarrow$) Os números \omterm{def:b3:galois:constructible}{construtíveis} formam um corpo: somas e diferenças
por paralelogramos (as paralelas são \omterm{def:b3:galois:constructible}{construtíveis}: baixe e levante
perpendiculares duas vezes --- a perpendicular clássica por um ponto
usa um círculo e dois arcos); produtos e quocientes pelas configurações
de Tales (dados os comprimentos $a, b$, construa $ab$ e
$a/b$ com triângulos semelhantes sobre duas semirretas). E o corpo
é fechado por raízes quadradas: para $a > 0$, o círculo de diâmetro
$1 + a$ e a perpendicular no ponto de junção encontram-se à altura
$\sqrt a$ (relação da média geométrica em um triângulo retângulo); para um
$w = \rho\eu^{\iu\theta}$ complexo, construa $\sqrt\rho$ e bissecte $\theta$ (a bissecção de
ângulo é uma construção com compasso). As partes real e imaginária dos
elementos de uma torre quadrática são, portanto, \omterm{def:b3:galois:constructible}{construtíveis}, por
indução sobre a torre: cada etapa adjunge raízes de uma quadrática,
exprimíveis por operações de corpo e uma raiz quadrada de um número já
construído (fórmula quadrática; em característica $0$).
\end{proof}

\begin{corollary}\label{cor:b3:galois:impossible}
Os três problemas clássicos são insolúveis com régua e
compasso:\index{régua e compasso}
\begin{enumerate}
  \item \emph{Duplicação do cubo}: $\sqrt[3]2$ tem \omterm{def:b3:galois:extension}{grau} $3$, que não é
        potência de $2$.
  \item \emph{Trissecção do ângulo}: trissectar $60^\circ$ exige $\cos 20^\circ$, raiz
        do \omterm{def:b3:rings:divisibility}{irredutível} $8X^3
        - 6X - 1$: \omterm{def:b3:galois:extension}{grau} $3$.
  \item \emph{Quadratura do círculo}: $\sqrt\pi$ é transcendente ($\pi$
        o é --- teorema de Lindemann, admitido aqui: sua demonstração
        pertence a um curso de teoria da transcendência).
\end{enumerate}
Além disso, o $n$-ágono regular é \omterm{def:b3:galois:constructible}{construtível} se, e somente se,
$\varphi(n)$ é uma potência de $2$ (Gauss--Wantzel; a implicação ``se''
usa o método do problema de fim de semana, e a implicação ``somente
se'' é o~\cref{thm:b3:galois:wantzel} aplicado a $\zeta_n$, de \omterm{def:b3:galois:extension}{grau} $\varphi(n)$). Para
$n = 7$: $\varphi(7) = 6$, e o heptágono regular é impossível; para
$n = 17$: $\varphi(17) = 16 = 2^4$, e ele é \omterm{def:b3:galois:constructible}{construtível} --- o problema de fim de
semana o constrói.
\end{corollary}

\begin{proof}
(1) $X^3 - 2$ é \omterm{def:b3:rings:divisibility}{irredutível} (\omterm{thm:b3:rings:criteria}{Eisenstein}). (2) De $\cos
3\theta = 4\cos^3\theta - 3\cos\theta$ com $3\theta =
60^\circ$: $8c^3 - 6c = 1$ para
$c = \cos 20^\circ$; a cúbica $8X^3 - 6X - 1$ não tem raiz racional (os candidatos $\pm1, \pm\frac
12, \pm\frac14, \pm\frac18$
falham) e, portanto, é \omterm{def:b3:rings:divisibility}{irredutível}: \omterm{def:b3:galois:extension}{grau} $3$. Um ângulo genérico
$60^\circ$ é \omterm{def:b3:galois:constructible}{construtível}, de modo que um trissector construiria $c$.
(3) Se $\sqrt\pi$ fosse \omterm{def:b3:galois:constructible}{construtível}, seria \omterm{def:b3:galois:algebraic}{algébrico} e, com ele, $\pi$.
Quanto ao enunciado sobre o $n$-ágono: o \omterm{def:b3:galois:extension}{grau} de $\zeta_n$ é $\varphi(n)$
(\cref{thm:b3:galois:cyclotomicirred}); a necessidade decorre de Wantzel; quanto à
suficiência, o \omterm{def:b3:galois:galois}{grupo de Galois}, abeliano de ordem $2^m$, admite uma
cadeia de subgrupos de índice $2$ (todo $2$-grupo finito admite:
\cref{exo:b3:groups:10}), cujos corpos fixos formam uma torre quadrática que
termina em $\Q(\zeta_n)$ (\cref{thm:b3:galois:fundamental}); conclua pelo~\cref{thm:b3:galois:wantzel}.
\end{proof}

\section{Solubilidade por radicais}

\begin{definition}\label{def:b3:galois:radical}
Uma extensão $L/K$ (com característica $0$ em toda esta seção)
é \emph{radical} se existe uma torre $K = F_0 \subseteq \dots
\subseteq F_r = L$ com $F_{i+1} = F_i(\alpha_i)$, $\alpha_i^{n_i} \in F_i$: cada etapa
adjunge uma raiz $n_i$-ésima. Um polinômio $P \in K[X]$ é \emph{solúvel
por radicais}\index{solúvel por radicais} se seu \omterm{thm:b3:galois:splitting}{corpo de decomposição}
está contido em alguma extensão radical de $K$.
\end{definition}

\begin{lemma}\label{lem:b3:galois:cyclicsteps}
Seja $K$ contendo uma raiz $n$-ésima primitiva da unidade
$\zeta$, isto é, $\zeta$ de ordem $n$ em $K^\times$, e seja
$a \in K^\times$. Então $K(\sqrt[n]a)/K$ é \omterm{def:b3:galois:galois}{de Galois} com grupo \emph{cíclico}. Reciprocamente
--- o que não será necessário adiante --- toda extensão cíclica de \omterm{def:b3:galois:extension}{grau}
$n$ é dessa forma. Além disso, $K(\zeta_n)/K$ é \omterm{def:b3:galois:galois}{de Galois} com grupo
\emph{abeliano}, para todo $K$ de característica $0$.
\end{lemma}

\begin{proof}
$X^n - a$ é \omterm{def:b3:galois:separable}{separável} ($\gcd$ com $nX^{n-1}$: $a \neq 0$) e se decompõe em $K(\alpha)$,
$\alpha^n = a$: suas raízes são os $\zeta^k\alpha \in K(\alpha)$. Logo $K(\alpha)/K$ é \omterm{def:b3:galois:galois}{de Galois}; a aplicação
$\sigma \mapsto \sigma(\alpha)/\alpha \in \mu_n = \langle
\zeta\rangle$ é um morfismo injetor ($\sigma\tau(\alpha) =
\sigma(\tau(\alpha)/\alpha \cdot \alpha) =
\tau(\alpha)/\alpha\cdot\sigma(\alpha)$, sendo o quociente um elemento de
$K$) em um grupo cíclico: $\operatorname{Gal}$ é cíclico. A recíproca é a teoria
de Kummer, de que não precisaremos (veja a observação abaixo). Quanto a
$K(\zeta_n)$: ela decompõe o \omterm{def:b3:galois:separable}{polinômio separável} $X^n -
1$, e $\sigma \mapsto a(\sigma)$ com $\sigma(\zeta_n) =
\zeta_n^{a(\sigma)}$
mergulha o grupo no abeliano $(\Z/n\Z)^\times$ como no~\cref{thm:b3:galois:cyclotomicirred} (a
injetividade só exige que $\zeta_n$ gere as raízes da unidade envolvidas).
\end{proof}

\begin{theorem}[Galois]\label{thm:b3:galois:solvable}
Sejam $K$ de característica $0$ e $P \in K[X]$ com \omterm{thm:b3:galois:splitting}{corpo de
decomposição} $L$. Se $P$ é \omterm{def:b3:galois:radical}{solúvel por radicais}, então $\operatorname{Gal}(L/K)$ é
um \omterm{def:b3:groups:derived}{grupo solúvel}. (A recíproca também é verdadeira; não precisaremos
dela.)\index{teorema de solubilidade de Galois}
\end{theorem}

\begin{proof}
Etapa 1: \emph{ampliar a torre \omterm{def:b3:galois:radical}{radical} até uma torre \omterm{def:b3:galois:galois}{de Galois}.} Seja
$L
\subseteq M$ com $M/K$ \omterm{def:b3:galois:radical}{radical}, de expoentes de raiz $n_1, \dots,
n_r$, e ponha $n = n_1\cdots n_r$.
Primeiro adjunte $\zeta_n$: a torre $K \subseteq K(\zeta_n) \subseteq M(\zeta_n)$ ainda é \omterm{def:b3:galois:radical}{radical} ($\zeta_n$ é uma raiz
da unidade: uma etapa \omterm{def:b3:galois:radical}{radical}, $\zeta_n^n = 1$), e suas etapas além da primeira
ocorrem sobre corpos que contêm as raízes da unidade necessárias. Em
seguida, substitua $M(\zeta_n)$ pelo composto $N$ de todos os $\sigma(M(\zeta_n))$, com
$\sigma$ percorrendo os (finitos) $K$-mergulhos de $M(\zeta_n)$ em um
\omterm{def:b3:galois:closure}{fecho algébrico} fixado: $N$ é o \omterm{thm:b3:galois:splitting}{corpo de decomposição} do produto
dos polinômios minimais de um conjunto gerador (característica
$0$: finito e \omterm{def:b3:galois:separable}{separável}), de modo que $N/K$ é \omterm{def:b3:galois:galois}{de Galois}; e
$N$ é \omterm{def:b3:galois:radical}{radical} sobre $K$: cada $\sigma(M(\zeta_n))$ é \omterm{def:b3:galois:radical}{radical} sobre $K$
(aplique $\sigma$ a uma torre \omterm{def:b3:galois:radical}{radical}), e um composto de extensões
radicais é \omterm{def:b3:galois:radical}{radical} (concatene as torres: se $F'/K$ é \omterm{def:b3:galois:radical}{radical} com
torre que adjunge $\beta_j$, as etapas do tipo $F''(\beta_j)$ permanecem radicais
sobre qualquer base maior).

Etapa 2: \emph{ler a \omterm{def:b3:groups:derived}{solubilidade} na torre \omterm{def:b3:galois:galois}{de Galois}.} Suponha, pois,
$L \subseteq N$, com $N/K$ \omterm{def:b3:galois:galois}{de Galois} e \omterm{def:b3:galois:radical}{radical}, de torre $K
\subseteq K(\zeta_n) = E_0 \subseteq E_1 \subseteq \dots
\subseteq E_s = N$, cada $E_{i+1} = E_i(\sqrt[n_i]{a_i})$
com $\zeta_{n_i} \in E_0 \subseteq E_i$. Sejam $G =
\operatorname{Gal}(N/K)$ e $G_i = \operatorname{Gal}(N/E_i)$: uma cadeia decrescente $G \supseteq G_0 \supseteq G_1 \supseteq
\dots \supseteq G_s = \{e\}$. Cada
$E_{i+1}/E_i$ é \omterm{def:b3:galois:galois}{de Galois} com grupo cíclico (\cref{lem:b3:galois:cyclicsteps}), de modo que, pelo
teorema fundamental aplicado à \omterm{def:b3:galois:galois}{extensão de Galois} $N/E_i$
(\cref{thm:b3:galois:fundamental}(3), com grupo ambiente $G_i$): $G_{i+1} \trianglelefteq G_i$ com $G_i/G_{i+1} \cong
\operatorname{Gal}(E_{i+1}/E_i)$
cíclico. Analogamente, $E_0/K$ é \omterm{def:b3:galois:galois}{de Galois} com grupo abeliano
$G/G_0$ (\cref{lem:b3:galois:cyclicsteps}). A cadeia exibe $G$ como \omterm{def:b3:groups:derived}{solúvel}
(\cref{prop:b3:groups:derived}). Por fim, $\operatorname{Gal}(L/K)$ é um \emph{quociente} de $G$:
$L/K$ é \omterm{def:b3:galois:galois}{de Galois} ($P$ é \omterm{def:b3:galois:separable}{separável} em característica $0$) e
a restrição $G
\to \operatorname{Gal}(L/K)$ é sobrejetora (\cref{thm:b3:galois:fundamental}(3) com $H =
\operatorname{Gal}(N/L)$); e
quocientes de \omterm{def:b3:groups:derived}{grupos solúveis} são \omterm{def:b3:groups:derived}{solúveis}.
\end{proof}

\begin{corollary}[Insolubilidade da quíntica]
\label{cor:b3:galois:quintic}
Há polinômios de \omterm{def:b3:galois:extension}{grau} $5$ sobre $\Q$ que não são \omterm{def:b3:galois:radical}{solúveis por
radicais}: por exemplo $X^5 - 4X + 2$, cujo \omterm{def:b3:galois:galois}{grupo de Galois} é $S_5$
(\cref{exo:b3:galois:11}), um grupo não \omterm{def:b3:groups:derived}{solúvel} (\cref{cor:b3:groups:snnotsolvable}). Nenhuma
fórmula geral por radicais pode existir para \omterm{def:b3:galois:extension}{grau}
$\geq 5$.\index{quíntica, insolubilidade}
\end{corollary}

\begin{remark}\label{rem:b3:galois:converse}
A recíproca do~\cref{thm:b3:galois:solvable} --- \omterm{def:b3:galois:galois}{grupo de Galois} \omterm{def:b3:groups:derived}{solúvel} implica
\omterm{def:b3:groups:derived}{solubilidade} por radicais --- demonstra-se descendo a \omterm{def:b3:groups:derived}{série derivada} e
mostrando que toda extensão cíclica (com raízes da unidade
suficientes) é \omterm{def:b3:galois:radical}{radical}, por meio dos \emph{resolventes de Lagrange};
ela explica por que os \omterm{def:b3:galois:extension}{graus} $2, 3, 4$ têm fórmulas: $S_2, S_3, S_4$ são \omterm{def:b3:groups:derived}{solúveis}
(\cref{ex:b3:groups:solvableexamples}). Deixamo-la admitida neste nível; um tratamento
completo pertence a um curso de mestrado, mas o~\cref{exo:b3:galois:8} a torna
concreta para o \omterm{def:b3:galois:extension}{grau} três.
\end{remark}

\section{Exercícios}

\begin{exercise}[$\star$]\label{exo:b3:galois:1}
Mostre que $[\Q(\sqrt2, \sqrt3):\Q] = 4$, que $\Q(\sqrt2 + \sqrt3) =
\Q(\sqrt2, \sqrt3)$, e calcule o \omterm{def:b3:galois:algebraic}{polinômio minimal} de $\sqrt2 + \sqrt3$
sobre $\Q$.
\end{exercise}

\begin{exercise}[$\star$]\label{exo:b3:galois:2}
Seja $\alpha = \sqrt[3]2$. Mostre que $\Q(\alpha)/\Q$ não é normal (exiba um mergulho $\Q(\alpha) \to \C$
cuja imagem não é $\Q(\alpha)$), determine o \omterm{thm:b3:galois:splitting}{corpo de decomposição} $L$ de
$X^3 -
2$ e $[L:\Q]$, e verifique que $\operatorname{Aut}_\Q(\Q(\alpha)) =
\{\mathrm{id}\}$: em extensões que não são
\omterm{def:b3:galois:galois}{de Galois}, o grupo de automorfismos pode ser muito menor do que o \omterm{def:b3:galois:extension}{grau}.
\end{exercise}

\begin{exercise}[$\star$]\label{exo:b3:galois:3}
Construa $\mathbb F_9$ como $\mathbb F_3[X]/(X^2+1)$ e encontre um gerador de $\mathbb F_9^\times$. Liste os
polinômios mônicos irredutíveis de \omterm{def:b3:galois:extension}{graus} $1, 2, 3$ sobre $\mathbb F_2$ e verifique
$X^8 - X = X(X+1)(X^3+X+1)(X^3+X^2+1)$ sobre $\mathbb F_2$.
\end{exercise}

\begin{exercise}[$\star\star$]\label{exo:b3:galois:4}
(a) Encontre todas as raízes primitivas módulo $7$ e módulo
$11$ (isto é, os geradores de $\mathbb F_7^\times$, $\mathbb F_{11}^\times$).
(b) Mostre que, para $p$ ímpar, $x \in \mathbb F_p^\times$ é um quadrado se, e somente
se, $x^{(p-1)/2} = 1$ (critério de Euler), e recupere o critério para $-1$
do~\cref{pb:b3:rings:1}.
\end{exercise}

\begin{exercise}[$\star\star$]\label{exo:b3:galois:5}
Mostre que $\mathbb F_{p^m} \cap \mathbb F_{p^n} = \mathbb
F_{p^{\gcd(m,n)}}$ e $\mathbb F_{p^m}\mathbb F_{p^n} = \mathbb
F_{p^{\operatorname{lcm}(m,n)}}$ dentro de um \omterm{def:b3:galois:closure}{fecho algébrico} fixado $\bar{\mathbb F}_p$, e
descreva $\operatorname{Gal}(\mathbb F_{p^n}/\mathbb F_{p^m})$ para $m
\mid n$.
\end{exercise}

\begin{exercise}[$\star\star$]\label{exo:b3:galois:6}
Seja $I_d(q)$ o número de polinômios mônicos irredutíveis de \omterm{def:b3:galois:extension}{grau}
$d$ sobre $\mathbb F_q$. Demonstre
\[
X^{q^n} - X \;=\; \prod_{d \mid n}\ \prod_{P \text{ irred.
monic, } \deg P = d} P ,
\qquad\text{logo}\qquad
q^n = \sum_{d \mid n} d\, I_d(q).
\]
Deduza $I_1, I_2, I_3, I_4$ explicitamente, e $I_d(q) \geq 1$ para todo $d$ (de modo que as
extensões $\mathbb F_{q^d}/\mathbb F_q$ existem como quocientes $\mathbb F_q[X]/(P)$ para todo $d$).
\end{exercise}

\begin{exercise}[$\star\star$]\label{exo:b3:galois:7}
Determine $\operatorname{Gal}(\Q(\sqrt2,\sqrt3)/\Q)$ e o reticulado completo dos corpos intermediários.
Mesma pergunta para o \omterm{thm:b3:galois:splitting}{corpo de decomposição} de $(X^2-2)(X^2-3)(X^2-6)$ --- o que você
observa?
\end{exercise}

\begin{exercise}[$\star\star$]\label{exo:b3:galois:8}
(A cúbica, resolvida por seu grupo) Seja $P = X^3 + pX + q \in
\Q[X]$ \omterm{def:b3:rings:divisibility}{irredutível} de raízes
$x_1, x_2, x_3$ e \omterm{thm:b3:galois:splitting}{corpo de decomposição} $L$. Sejam $\delta = (x_1 - x_2)(x_1 - x_3)(x_2 - x_3)$ e $\Delta = \delta^2 = -4p^3 - 27q^2$ (admita essa
identidade clássica ou verifique-a expandindo funções simétricas).
(a) Mostre que $\operatorname{Gal}(L/\Q) \cong A_3$ ou $S_3$, conforme $\Delta$ seja ou não um
quadrado em $\Q$.
(b) Com $j = \zeta_3$, defina os resolventes de Lagrange $u = x_1
+ jx_2 + j^2x_3$ e $v = x_1 + j^2x_2 + jx_3$. Mostre
que $u^3 + v^3
= -27q$ e $uv = -3p$, e resolva para $u^3, v^3$: as fórmulas de Cardano caem
sozinhas. Onde foi usada a \omterm{def:b3:groups:derived}{solubilidade} de $S_3$?
\end{exercise}

\begin{exercise}[$\star\star$]\label{exo:b3:galois:9}
Em $\Q(\zeta_5)$: mostre que o único subcorpo quadrático é $\Q(\sqrt5)$, por meio das
somas de Gauss $\eta_0 = \zeta_5 + \zeta_5^4$, $\eta_1 = \zeta_5^2 + \zeta_5^3$: calcule $\eta_0 + \eta_1$ e $\eta_0\eta_1$, e deduza $\cos\frac{2\pi}5 = \frac{\sqrt5 -
1}4$.
Conclua que o pentágono regular é \omterm{def:b3:galois:constructible}{construtível}.
\end{exercise}

\begin{exercise}[$\star\star\star$]\label{exo:b3:galois:10}
Sejam $K = \mathbb F_p(S, T)$ (funções racionais em duas indeterminadas) e $L = K(S^{1/p}, T^{1/p})$.
(a) Mostre que $[L:K] = p^2$ e que $\alpha^p \in K$ para todo $\alpha \in L$.
(b) Deduza que $L/K$ \emph{não} é simples: não existe elemento
primitivo --- a inseparabilidade é fatal para o~\cref{thm:b3:galois:primitive}.
\end{exercise}

\begin{exercise}[$\star\star\star$]\label{exo:b3:galois:11}
Sejam $P = X^5 - 4X + 2$ e $G$ seu \omterm{def:b3:galois:galois}{grupo de Galois} sobre $\Q$, agindo sobre
as $5$ raízes.
(a) Mostre que $P$ é \omterm{def:b3:rings:divisibility}{irredutível} e deduza $5 \mid \abs G$; conclua que
$G$ contém um $5$-ciclo (Cauchy, \cref{thm:b3:groups:cauchy}).
(b) Mostre, estudando as variações de $x \mapsto x^5 - 4x +
2$, que $P$ tem exatamente
$3$ raízes reais; deduza que a conjugação complexa se restringe a
uma transposição em $G$.
(c) Mostre que um subgrupo de $S_5$ que contém uma transposição e
um $5$-ciclo é $S_5$ \emph{(conjugue a transposição por
potências do ciclo)}. Conclua $G \cong S_5$ e, com o~\cref{thm:b3:galois:solvable}, que
$P$ não é \omterm{def:b3:galois:radical}{solúvel por radicais}.
\end{exercise}

\begin{exercise}[$\star\star\star$]\label{exo:b3:galois:12}
(A quártica diedral) Sejam $\alpha = \sqrt[4]2$ e $L =
\Q(\alpha, \iu)$, o \omterm{thm:b3:galois:splitting}{corpo de decomposição} de
$X^4 - 2$ sobre $\Q$.
(a) Mostre que $[L : \Q] = 8$ e que $G =
\operatorname{Gal}(L/\Q)$ é gerado por $\sigma\colon
\alpha \mapsto \iu\alpha,\ \iu \mapsto \iu$ e pela conjugação
complexa $\tau$, com $\sigma^4 = \tau^2 = e$ e $\tau\sigma\tau = \sigma^{-1}$: $G \cong D_4$.
(b) Liste o reticulado de subgrupos de $D_4$ (dez subgrupos) e
associe cada um ao seu corpo fixo; verifique em particular que $\Q(\sqrt2)$,
$\Q(\iu)$, $\Q(\iu\sqrt2)$ são os três subcorpos quadráticos, e localize $\Q(\alpha)$,
$\Q(\iu\alpha)$, $\Q(\sqrt2, \iu)$.
(c) Quais corpos intermediários são \omterm{def:b3:galois:galois}{de Galois} sobre $\Q$? Confronte
sua resposta com os \omterm{def:b3:groups:normal}{subgrupos normais} de $D_4$ e explique por que
$\Q(\alpha)/\Q$ falha enquanto $\Q(\sqrt2)/\Q$ funciona.
\end{exercise}

\section{Problema: Gauss e o 17-ágono regular}

\begin{problem}[{Problema de fim de semana --- construtibilidade do
17-ágono}]\label{pb:b3:galois:1}
Em 30 de março de 1796, Gauss, aos dezenove anos, mostrou que o
$17$-ágono regular é \omterm{def:b3:galois:constructible}{construtível} --- o primeiro avanço sobre a
questão desde a Antiguidade. Reconstruímos seu cálculo com as
ferramentas deste capítulo. Ponha $\zeta = \eu^{2\iu\pi/17}$, $L =
\Q(\zeta)$, $G = \operatorname{Gal}(L/\Q)$.

\medskip
\noindent\textbf{Parte I --- O grupo e sua filtração.}
\begin{enumerate}
  \item Justifique: $[L:\Q] = 16$, $G \cong (\Z/17\Z)^\times$, cíclico de ordem $16$.
        Verifique que $3$ é gerador de $(\Z/17\Z)^\times$ \emph{(calcule as
        potências de $3$ módulo $17$: $3, 9, 10, 13, 5, 15, 11, 16, \dots$)}.
  \item Sejam $\sigma \in G$ com $\sigma(\zeta) = \zeta^3$, e $H_k = \langle \sigma^{2^k}\rangle$ para $k = 0, \dots,
        4$. Mostre que $G = H_0 \supset H_1 \supset H_2 \supset
        H_3 \supset H_4 = \{e\}$
        com cada índice $[H_k :
        H_{k+1}] = 2$, e que os corpos fixos $\Q = L_0
        \subset L_1 \subset L_2 \subset L_3 \subset L_4 = L$ formam uma
        torre de extensões quadráticas.
  \item Conclua \emph{a priori}, usando o~\cref{thm:b3:galois:wantzel}, que
        $\zeta$ --- e portanto o $17$-ágono --- é \omterm{def:b3:galois:constructible}{construtível}.
        O restante do problema torna a torre explícita.
\end{enumerate}

\medskip
\noindent\textbf{Parte II --- Os períodos de comprimento 8.}
Defina os \emph{períodos de Gauss}
\[
\eta_0 = \sum_{k \text{ par}} \zeta^{3^k \bmod 17}
= \zeta^{1} + \zeta^{9} + \zeta^{13} + \zeta^{15} + \zeta^{16} +
\zeta^{8} + \zeta^{4} + \zeta^{2},
\qquad
\eta_1 = \sum_{k \text{ ímpar}} \zeta^{3^k \bmod 17}.
\]
\begin{enumerate}[resume]
  \item Mostre que $\eta_0, \eta_1$ são fixados por $H_1$ e trocados por
        $\sigma$; deduza $\eta_0, \eta_1 \in L_1$ e que eles são as duas raízes de
        uma quadrática sobre $\Q$.
  \item Calcule $\eta_0 + \eta_1 = -1$. Mostre que $\eta_0\eta_1 = -4$ \emph{(cada produto $\zeta^a\zeta^b$ é
        algum $\zeta^c$, $c \neq 0$; conte quantas vezes cada $c$ ocorre,
        ou argumente que o produto é um inteiro racional fixado por
        $G$, igual à soma sobre todos os $64$ produtos, e use
        que cada resto não nulo aparece o mesmo número de vezes)}.
  \item Deduza $\eta_0 = \frac{-1 + \sqrt{17}}2$, $\eta_1 =
        \frac{-1-\sqrt{17}}2$ (identifique numericamente qual é qual:
        $\eta_0 \approx 1.56$) e $L_1 =
        \Q(\sqrt{17})$.
\end{enumerate}

\medskip
\noindent\textbf{Parte III --- Períodos de comprimento 4 e 2.}
Defina
\[
\beta_0 = \zeta + \zeta^{13} + \zeta^{16} + \zeta^{4},\quad
\beta_1 = \zeta^3 + \zeta^5 + \zeta^{14} + \zeta^{12},\quad
\beta_2 = \zeta^9 + \zeta^{15} + \zeta^{8} + \zeta^{2},\quad
\beta_3 = \zeta^{10} + \zeta^{11} + \zeta^{7} + \zeta^{6}.
\]
\begin{enumerate}[resume]
  \item Mostre que $\beta_0 + \beta_2 = \eta_0$, $\beta_1 + \beta_3 =
        \eta_1$, e que $\beta_0, \beta_2$ são fixados por
        $H_2$ e trocados por $\sigma^2$.
  \item Calcule $\beta_0\beta_2 = -1$ e $\beta_1\beta_3 = -1$ \emph{(expanda: os dezesseis expoentes
        obtidos cobrem $1, \dots, 16$ exatamente uma vez)}.
  \item Deduza $\beta_0 = \frac{\eta_0 + \sqrt{\eta_0^2 + 4}}2$ (confira o sinal numericamente: $\beta_0 \approx 2.05$) e a
        fórmula análoga para $\beta_1$; logo $L_2 =
        \Q(\beta_0)$, quadrática sobre
        $L_1$.
  \item Sejam $\gamma_0 = \zeta + \zeta^{16} =
        2\cos\frac{2\pi}{17}$ e $\gamma_1 = \zeta^{13} +
        \zeta^4$. Mostre que $\gamma_0 + \gamma_1 = \beta_0$ e $\gamma_0\gamma_1 = \beta_1$, de modo que
        $\gamma_0 =
        \frac{\beta_0 + \sqrt{\beta_0^2 - 4\beta_1}}2$.
  \item Monte a cadeia de fórmulas que exprime $\cos\frac{2\pi}{17}$ por radicais
        quadráticos encaixados e dê uma verificação decimal ($\cos\frac{2\pi}{17} \approx 0.93247$).
\end{enumerate}

\medskip
\noindent\textbf{Parte IV --- Epílogo.}
\begin{enumerate}[resume]
  \item Onde exatamente o argumento usou que $17$ é um \emph{primo
        de Fermat} ($17 = 2^{2^2} + 1$)? Mostre que, para um primo $p$, o
        $p$-ágono regular é \omterm{def:b3:galois:constructible}{construtível} se, e somente se, $p
        = 2^{2^t} + 1$
        para algum $t$ \emph{(se $p - 1 = 2^m$, mostre que $m$ tem de
        ser ele próprio uma potência de $2$)}.
  \item Deduza a lista completa dos $n$-ágonos regulares
        \omterm{def:b3:galois:constructible}{construtíveis} para $n \leq 20$, usando o critério de Gauss--Wantzel
        do~\cref{cor:b3:galois:impossible}.
\end{enumerate}

\medskip
\noindent\textbf{Parte V --- Somas de Gauss e reciprocidade
quadrática.} Os períodos da Parte II escondem um tesouro. Para um primo
ímpar $p$, o \emph{símbolo de Legendre} $\bigl(\frac
ap\bigr)$ vale $+1$ se
$a$ é um quadrado não nulo módulo $p$, $-1$ se não é, e
$0$ se $p \mid a$; o~\cref{exo:b3:galois:4}(b) (critério de Euler) dá $\bigl(\frac ap\bigr) \equiv
a^{(p-1)/2} \pmod p$,
donde a multiplicatividade. Escreva $\zeta =
\eu^{2\iu\pi/p}$, $p^* = (-1)^{(p-1)/2}p$, e defina a
\emph{soma de Gauss}
\[
g \;=\; \sum_{a=1}^{p-1}\Bigl(\frac ap\Bigr)\zeta^a .
\]
\begin{enumerate}[resume]
  \item Mostre que $\sum_{a=1}^{p-1}\bigl(\frac ap\bigr) = 0$ (há tantos quadrados quanto não quadrados) e
        demonstre a forma alternativa $g = \sum_{a=0}^{p-1}\zeta^{a^2}$ \emph{(cada quadrado não
        nulo é atingido duas vezes, e $\sum_{a}\zeta^a = 0$)}. Para $p = 17$:
        relacione $g$ com os períodos da Parte II --- mostre que
        $g = \eta_0 - \eta_1$ \emph{(os quadrados módulo $17$ são exatamente as
        potências pares do gerador $3$)}.
  \item Demonstre que $g^2 = p^*$: expanda
        \[
        g^2 = \sum_{a,b\neq0}\Bigl(\frac{ab}p\Bigr)
        \zeta^{a+b}
        = \sum_{c}\ \sum_{a \neq 0}\Bigl(\frac{a(c -
        a)}p\Bigr)\zeta^{c}
        \]
        (ponha $b = c - a$), substitua $c - a = at$ para avaliar a soma interna
        como $\bigl(\frac{-1}p\bigr)(p - 1)$ para $c = 0$ e $-\bigl(\frac{-1}p\bigr)$ nos demais casos, e
        conclua com a questão 14. Confira numericamente: para
        $p =
        17$, $(\eta_0 - \eta_1)^2 = 17$ (Parte II).
  \item Deduza $\sqrt{p^*} \in \Q(\zeta_p)$ e conclua que o \emph{único} subcorpo quadrático
        de $\Q(\zeta_p)$ é $\Q(\sqrt{p^*})$ --- único porque $\operatorname{Gal}(\Q(\zeta_p)/\Q)$ é cíclico
        (\cref{thm:b3:galois:cyclotomicirred}) e um grupo cíclico tem exatamente um subgrupo
        de índice $2$. (Todo corpo quadrático mergulha em algum
        corpo ciclotômico --- este é o primeiro caso do teorema de
        Kronecker--Weber, cuja forma geral está muito à frente.)
  \item Seja agora $q \neq p$ outro primo ímpar. Trabalhando no anel $\Z[\zeta]$
        módulo $q$, demonstre
        \[
        g^q \equiv \Bigl(\frac qp\Bigr)\,g \pmod{q\Z[\zeta]}
        \]
        \emph{(sonho do calouro: $(x + y)^q \equiv x^q + y^q$ módulo $q$ em qualquer anel
        comutativo; depois $g^q \equiv
        \sum_a\bigl(\frac ap\bigr)^q\zeta^{aq}$, reindexe $b =
        aq$ e ponha
        $\bigl(\frac{q^{-1}}p\bigr) =
        \bigl(\frac qp\bigr)$ em evidência)}.
  \item Por outro lado, $g^q = g\,(g^2)^{(q-1)/2} =
        g\,(p^*)^{(q-1)/2}$; usando o critério de Euler módulo
        $q$, deduza $g^q \equiv \bigl(\frac{p^*}q\bigr)g
        \pmod{q\Z[\zeta]}$ e, então --- multiplicando as duas
        expressões de $g^q$ por $g$ e usando $g^2 = p^*$, inversível
        módulo $q$ ---, conclua
        \[
        \Bigl(\frac qp\Bigr) = \Bigl(\frac{p^*}q\Bigr) .
        \]
        \emph{(Por que uma congruência entre os inteiros $\pm p^*$ módulo
        $q\Z[\zeta]$ implica sua igualdade? Intersecte com $\Z$.)}
  \item Desenvolva $\bigl(\frac{p^*}q\bigr) =
        \bigl(\frac{-1}q\bigr)^{(p-1)/2}\bigl(\frac pq\bigr)$ e $\bigl(\frac{-1}q\bigr) = (-1)^{(q-1)/2}$ para obter a \emph{lei da
        reciprocidade quadrática}:
        \[
        \Bigl(\frac pq\Bigr)\Bigl(\frac qp\Bigr)
        = (-1)^{\frac{p-1}2\cdot\frac{q-1}2} .
        \]
        Verifique-a em $(p, q) = (17, 3)$, listando os quadrados módulo $17$ e
        módulo $3$, e use-a para decidir em três linhas se $x^2 \equiv 219 \pmod{383}$
        tem solução ($383$ é primo, $219 = 3\cdot73$).
\end{enumerate}

\medskip
\noindent\textbf{Parte VI --- Contando polinômios irredutíveis: o
teorema dos números primos de $\mathbb F_q[X]$.} Fixe uma potência de primo
$q$ e seja $N_q(n)$ o número de polinômios \emph{mônicos
irredutíveis} de \omterm{def:b3:galois:extension}{grau} $n$ sobre $\mathbb F_q$; recorde do~\cref{exo:b3:galois:6}
a fatoração de $X^{q^n} - X$ e a identidade $q^n = \sum_{d\mid n}d\,N_q(d)$, que agora invertemos,
reinterpretamos e exploramos.
\begin{enumerate}[resume]
  \item (Palavras) Diga que uma palavra $w \in \mathbb F_q^n$ é \emph{primitiva} se
        não é uma potência $u^{n/d} =
        u\cdots u$ de uma palavra $u$ estritamente
        mais curta, e seja $A(d)$ o número de palavras primitivas
        de comprimento $d$. Mostre que toda palavra de comprimento
        $n$ é, de maneira única, uma potência de uma palavra
        primitiva de algum comprimento $d \mid n$, de modo que $q^n = \sum_{d \mid n}A(d)$;
        comparando com o~\cref{exo:b3:galois:6}, conclua $A(d) = d\,N_q(d)$ para todo
        $d$ e explique essa coincidência por uma bijeção
        explícita: um elemento $\alpha \in
        \mathbb F_{q^d}$ de \omterm{def:b3:galois:extension}{grau} $d$ tem \omterm{def:b3:groups:action}{órbita} de
        \omterm{thm:b3:galois:finitefields}{Frobenius} $(\alpha, \alpha^q, \dots, \alpha^{q^{d-1}})$ com exatamente $d$ elementos distintos, e
        os elementos de \omterm{def:b3:galois:extension}{grau} $d$ correspondem $d$ a um aos
        irredutíveis de \omterm{def:b3:galois:extension}{grau} $d$.
  \item Demonstre a \emph{fórmula de inversão de Möbius}: se $f(n) =
        \sum_{d\mid n}g(d)$
        para todo $n$, então $g(n) =
        \sum_{d\mid n}\mu(d)\,f(n/d)$, onde $\mu$ é a função de
        Möbius ($\mu(m) = (-1)^{\#\text{prime
        factors}}$ se $m$ é livre de quadrados, $0$ caso
        contrário) \emph{(lema-chave: $\sum_{d \mid m}\mu(d) = 0$ para $m >
        1$ ---
        emparelhe os divisores com e sem um fator primo fixado)}.
        Deduza
        \[
        N_q(n) = \frac1n\sum_{d \mid n}\mu(d)\,q^{n/d} .
        \]
  \item Mostre que $N_q(n) \geq \frac1n\bigl(q^n - 2q^{n/2}\bigr) > 0$ para todo $n \geq 1$: uma nova demonstração de que
        $\mathbb F_{q^n}$ existe para todo $n$. Interprete o termo dominante:
        um polinômio mônico aleatório de \omterm{def:b3:galois:extension}{grau} $n$ é \omterm{def:b3:rings:divisibility}{irredutível} com
        probabilidade $\sim \frac1n$ --- o análogo \omterm{prop:b3:galois:perfect}{perfeito} do teorema dos
        números primos, com $\log x$ trocado por $n$; verifique
        numericamente para $q = 2$, $n \leq 4$ (o~\cref{exo:b3:galois:6} dá as
        contagens).
  \item Demonstre a companheira multiplicativa da questão 21:
        \[
        \prod_{\substack{\pi \text{ mônico irred.}\\ \deg\pi =
        n}}\pi \;=\; \prod_{d \mid n}
        \bigl(X^{q^d} - X\bigr)^{\mu(n/d)}
        \]
        \emph{(inversão de Möbius no grupo abeliano das funções
        racionais não nulas)}; verifique-a à mão para $q
        = 2$,
        $n = 2$: $(X^4 - X)/(X^2 - X) = X^2 + X + 1$.
\end{enumerate}

\medskip
\noindent\textbf{Parte VII --- Duas codas.}
\begin{enumerate}[resume]
  \item (O segundo complemento) O método da Parte V também calcula
        $\bigl(\frac2q\bigr)$. Sejam $\omega = \eu^{2\iu\pi/8}$ e $g = \omega + \omega^{-1}$. Mostre que $g^2 = 2$ \emph{($\omega^2 = \iu$)};
        depois, para um primo ímpar $q$, demonstre em $\Z[\omega]$ módulo
        $q$ que
        \[
        g^q \equiv \omega^q + \omega^{-q} \pmod{q\Z[\omega]},
        \]
        e que o lado direito é igual a $g$ se $q \equiv \pm1
        \pmod 8$ e a $-g$
        se $q \equiv \pm3 \pmod 8$. Comparando com $g^q = g\,(g^2)^{(q-1)/2} \equiv
        \bigl(\frac2q\bigr)g$ como na questão 18, conclua
        \[
        \Bigl(\frac2q\Bigr) = (-1)^{(q^2-1)/8},
        \]
        verificando que $(q^2 - 1)/8$ é par exatamente quando $q
        \equiv \pm1 \pmod 8$. Confira:
        $2$ é quadrado módulo $7$ e módulo $17$
        ($3^2$ e $6^2$), e não é módulo $3$ nem módulo
        $5$.
  \item (A função zeta de $\mathbb F_q[X]$) Demonstre a identidade de séries
        formais de potências em $t$:
        \[
        \prod_{n \geq 1}\bigl(1 - t^n\bigr)^{-N_q(n)}
        = \frac1{1 - qt}
        \]
        \emph{(fatoração única em mônicos irredutíveis: expanda cada
        fator como série geométrica e conte os polinômios mônicos de
        \omterm{def:b3:galois:extension}{grau} $n$)}. Recupere a identidade $q^m = \sum_{d \mid m}d\,N_q(d)$ do~\cref{exo:b3:galois:6}
        tomando logaritmos. Confira à mão o coeficiente de $t^2$
        para $q = 2$, e use a fórmula da questão 21 para calcular
        $N_2(6) = 9$, verificando $2^6 = 1\cdot2 + 2\cdot1 + 3\cdot2 +
        6\cdot9$.
\end{enumerate}
\end{problem}
